\paragraph{Theorem F.2.16 (Proved: an unbounded integral family).}
There are infinitely many simultaneous rational conjugacy classes
of integral rank-six tuples of the linear type $(5,4)$ in F.2,
with primitive maximal square-zero cusp.  An explicit infinite
subfamily has free affine lifts of orders five and four.
Apart from the parameter zero, the rank-five irreducible quotient
in the family is not self-dual.  Thus no finite bound on rational
conjugacy classes is valid for the unrestricted input tuples.
This does not assert infinitely many three-connected survivors.

\begin{proof}
All the following expressions are explicit.  Let $v\in\mathbb Z$,
$q=v^2-v+1$, and let $A$ be the permutation matrix
$Ae_j=e_{j+1}$ on $\mathbb Z^5$, with subscripts taken modulo five.
Put
\[
 X=(e_0,\ e_2+ve_4),\qquad
 Y=\begin{pmatrix}
 0&-q&-v^2&-1&v\\
 0&1-vq&-vq&v^2-v&q
 \end{pmatrix},\qquad N=XY,
 \tag{F.U1}
\]
and set $C=(I+N)A$, $B=C^{-1}=A^{-1}(I-N)$.
Direct multiplication gives
\[
 YX=0,\quad Y(A+I)^{-1}X=I_2,\quad
 A^5=C^4=I_5,\quad
 \chi_C(x)=(x-1)(x+1)^2(x^2+1).
 \tag{F.U2}
\]
The image of $X$ is primitive, and $Y$ is onto $\mathbb Z^2$.
For the latter assertion, $2(A+I)^{-1}X$ is integral, so the
cokernel of $Y$ is killed by two; reduction of (F.U1) in the
two cases $v\equiv0,1\pmod2$ has rank two.  Thus $N$ has
rank two, $N^2=0$, and its two nonzero Smith factors are one.

The common invariant space of $A,B$ is zero: the only
$A$-fixed line is generated by $\mathbf1=(1,1,1,1,1)^t$,
and the first entry of $Y\mathbf1$ is $-2q\ne0$.
There are no common invariant covectors either, since
$\mathbf1^tX=(1,v+1)$ and $Y$ has rank two.
The characteristic polynomial of $A$ is $(x-1)\Phi_5(x)$.
A proper rational invariant subspace must therefore have
dimension one or four.  A one-dimensional subrepresentation
or quotient has $A$ acting as one; the unipotent product then
forces $B$ also to act as one.  The two vanishing assertions
exclude both possibilities.  The rank-five pair is irreducible.

Here is an explicit primitive rank-six extension, including
the integral lattice.  Let $P$ have columns
$\mathbf1/q,e_1,e_2,e_3,e_4$, and put
$A'=P^{-1}AP$, $B'=P^{-1}BP$.  These matrices are integral:
$A\mathbf1=\mathbf1$ and
\[
 Y\mathbf1=2q\binom{-1}{1-v}
\]
show that the index-$q$ overlattice is stable.  Define
\[
 s=v^2+v-1,\qquad r=-\frac{(v+1)(v^2-2)}2,
 \qquad
 \beta=(s,\ s-vr,\ v^2s+vr,\ r,\ -vs-r).
 \tag{F.U3}
\]
The number $r$ is integral for every integer $v$.
Set
\[
 \widehat A=\begin{pmatrix}1&0\\0&A'\end{pmatrix},
 \qquad
 \widehat B=\begin{pmatrix}1&\beta\\0&B'\end{pmatrix}.
 \tag{F.U4}
\]
They have orders five and four, and
$(\widehat A\widehat B-I)^2=0$.
To check primitivity without a parameter bound, write
$X'=P^{-1}X$, $Y'=YP$.  Two rows of $X'$ have determinant
$-1$.  The quotient of $\mathbb Z^5$ by its image has
coordinates indexed by $0,3,4$, and the induced matrix $Y'$
there is
\[
 \begin{pmatrix}-2&-1&v\\2(1-v)&v^2-v&q\end{pmatrix}.
\]
Its kernel has primitive generator
$(-v^3+v-1,\,2,\,2(1-v^2))^t$: its first coordinate is
odd and the second is two.  Lift this to
\[
 z=(-v^3+v-1,0,0,2,2(1-v^2))^t.
\]
The same three entries are its signed two-by-two minors,
whose gcd is one; hence $Y'$ is onto $\mathbb Z^2$.
Then $Y'z=0$ and
\[
 \beta z=s(v^3-v-1)+2v^2r=1.
 \tag{F.U5}
\]
Thus the rank-two primitive cusp in the quotient acquires a
primitive third image direction in the extension.  All three
nonzero Smith factors of $\widehat A\widehat B-I$ are one.

For later homology comparisons use the following constant
change of lattice:
\[
 L=\frac14\begin{pmatrix}
 -2&3&-5&-1&5&-3\\
 0&2&-2&2&-2&2\\
 0&-4&0&0&4&-4\\
 0&0&0&-4&4&-4\\
 0&0&4&-4&4&-4\\
 0&0&0&-4&4&0
 \end{pmatrix},\quad
 A_v=L^{-1}\widehat A L,\quad B_v=L^{-1}\widehat B L.
 \tag{F.U6}
\]
Here $\det L=-1/4$, $L^{-1}$ is integral, and all entries of
$A_v,B_v$ are polynomials in $v$ with integer coefficients.
The primitive-cusp assertion is unchanged at odd primes;
at two, direct reduction for $v=0,1$ gives rank three.
At $v=0$, (F.U6) is exactly the pair in Theorem F.2.5.

For the infinite subfamily $v\in20\mathbb Z$, take homology
monodromies $M_1=A_v^t$, $M_2=B_v^t$ and
$M_0=(M_1M_2)^{-1}$.  The reversed product has the same
primitive square-zero cusp type.  Write
$a=(1,3,-2,-1,2,0)^t$ and $e=(1,0,0,0,0,0)^t$, and set
\[
 w_1=\frac15\sum_{j=0}^4(A_v^t)^j a,
 \qquad
 w_2=-\frac14\sum_{j=0}^3(B_v^t)^j e.
 \tag{F.U7}
\]
These vectors are integral.  Indeed, the polynomial matrices
reduce to their parameter-zero values modulo five and four,
respectively, and $a,e$ are fixed by the corresponding
parameter-zero matrices.  They are fixed by $M_1,M_2$, with
first coordinates $1,-1$.  Both monodromies preserve the
first-coordinate covector $\ell$.  Thus the affine actions
$x\mapsto M_i x+w_i/m_i$, with $(m_1,m_2)=(5,4)$, are free:
every nontrivial power translates the $\ell$-circle nontrivially.

Finally, let $H=NA|_{\operatorname{im}N}$ and
$H'=NA^{-1}|_{\operatorname{im}N}$.  The construction gives
\[
 \det H=1-v^2q,\quad \det H'=1+v^2q,
 \qquad\operatorname{tr}(A^2C^2)=v^5-v^4+3.
 \tag{F.U8}
\]
These are simultaneous-conjugacy invariants; the last is
unbounded on the positive multiples of twenty.  If the
irreducible rank-five representation were self-dual, its
invariant bilinear form would have a nonzero symmetric part,
which is nondegenerate by irreducibility; an alternating
nondegenerate form is impossible in odd dimension.  Relative
to that symmetric form, $A,C$ are orthogonal and
$N^*=-N$, so $I+tN$ is orthogonal for every $t$.
The odd-dimensional characteristic polynomial of $(I+tN)A$
is then anti-reciprocal.  Its quadratic coefficient in $t$ is
\[
 x^2\bigl((\det H)x-\det H'\bigr),
\]
by the determinant lemma, so anti-reciprocity forces
$\det H=\det H'$.  Formula (F.U8) contradicts this whenever
$v\ne0$.

All the matrix identities, coefficient integrality, both
mod-two checks, and the affine formula are checked symbolically
by \texttt{code/fourfolds\_unbounded\_family.py}.  Its separate
Smith diagnostics over $-6\le v\le10$ are explicitly bounded
and are not the basis for the universal assertions above.
\end{proof}

\paragraph{Proposition F.2.17 (Proved: an odd torsion obstruction in this lattice).}
For the integral tuple $(A_v,B_v)$ in (F.U6), any simply
connected filling with free finite affine caps and the fixed
primitive Kuhn cusp has a subgroup
\[
 \mathbb Z/(v^2-v+1)\ \subset\operatorname{Tor}H_2(M;\mathbb Z).
 \tag{F.U9}
\]
In particular, in this specified family of lattices only
$v=0$ and $v=1$ can be three-connected, regardless of the
affine multiplicities and twists.  The stronger exclusion for every stable lattice in each
rational class is proved in Theorem F.2.19 below.

\begin{proof}
First compute the global coinvariants in the hatted lattice.
In the rank-five quotient, the $A'$ relations identify
$f_1=f_2=f_3=f_4=t$ and $qf_0=5t$.
The number $q=v^2-v+1$ is relatively prime to five, as its
five residue values show.  Thus these coinvariants are
infinite cyclic, with $f_0=5h$ and $t=qh$.
The $B'$ relations become $q\kappa_j h$, where
\[
 (\kappa_0,\ldots,\kappa_4)=
 (2v^2,\ v(v^3+v-1),\ v(v^3+v+1),
       -v^3+v+1,\ -v^3-v-1).
 \tag{F.U10}
\]
Their gcd is one: it divides $2v$ by subtracting the second
and third entries, and it divides $v^3-v-1$ by the fourth;
the latter is odd and is $-1$ modulo $v$.
Hence the rank-five global coinvariants are $\mathbb Z/q$.
The extension line is killed in rank-six coinvariants by
(F.U5), since
$(\widehat A\widehat B-I)(0,z)^t=e$.
Consequently the hatted rank-six coinvariants are also
$\mathbb Z/q$.  The lattice change (F.U6) is an isomorphism
over every $\mathbb Z_{(p)}$ with $p$ odd, so the odd part of
the actual rank-six global coinvariants is precisely
$\mathbb Z/q$.

Let $V$ be the degree-one fibre cohomology local system.
The only difference between $R^1g_*\mathbb Z$ and $j_*V$
is a finite cokernel at finite caps; the cusp degree-one
specialization is onto.  Thus
\[
 H^2(\mathbb P^1,R^1g_*\mathbb Z)
 =H^2(\mathbb P^1,j_*V)=V_\Gamma,
\]
including its $q$-primary part just computed.  The space of
global degree-two invariants is zero: the irreducible odd
rank-five quotient has no invariant alternating tensor,
and the extension filtration
$0\to W\to\bigwedge^2V\to\bigwedge^2W\to0$
gives $(\bigwedge^2V)^\Gamma=0$.

It remains to check every possible source of the differential
$d_2:H^0(\mathbb P^1,R^2g_*\mathbb Z)
\to H^2(\mathbb P^1,R^1g_*\mathbb Z)$.
The cusp has no kernel in degree two.  At a finite cap the
kernel is torsion in $H^2$ of that cap, namely torsion in its
first homology.  Simple connectivity of the total filling
forces the distinguished affine numerator to be prime to the
actual affine multiplicity, by the determinant argument of
F.2.12.  The cap is therefore a mapping torus of a torus
automorphism whose order divides the linear order, five or
four.  Its first-homology torsion is supported only at the
primes five or two, respectively: averaging over the finite
order kills the torsion of the vertical coinvariants.
Hence the source of this $d_2$ has no torsion at any prime
dividing $q$, since $q$ is odd and prime to five.

The subgroup $\mathbb Z/q$ in bidegree $(2,1)$ survives.
There are no longer differentials on a two-dimensional base,
and this terminal filtration term is a subgroup of
$H^3(M;\mathbb Z)$.  Universal coefficients identify its
torsion with a subgroup of $\operatorname{Tor}H_2(M;\mathbb Z)$.
This proves (F.U9).  Finally, $v^2-v+1=1$ for an integer $v$
exactly when $v=0$ or $v=1$.
\end{proof}

\paragraph{Theorem F.2.18 (Proved: a survivor with a non-self-dual quotient).}
The parameter $v=1$ in (F.U6), with affine orders five and four
and twists
\[
 w_1=(1,3,-2,-1,2,0)^t,\qquad
 w_2=(-1,0,0,0,0,0)^t,
\]
gives a three-connected smooth eight-manifold with
$H_4=\mathbb Z^4$ and all other reduced homology zero except
$H_8=\mathbb Z$.  It is almost complex and is oriented
diffeomorphic to $3\mathbb{HP}^2\#(-\mathbb{HP}^2)$.
Its irreducible rank-five quotient has no invariant bilinear
form, so integral three-connectivity does not force that
quotient to be orthogonal.  It has bounded-cusp holomorphic
periods by Theorem F.3.11; simultaneous finite holomorphic
extensions are not asserted here.

\begin{proof}
The first cohomological matrix is the same $A$ as in
Theorem F.2.5.  The second is
\[
 B_1=\begin{pmatrix}
 1&0&0&0&0&0\\
 0&0&-1&0&1&-1\\
 0&0&-1&0&0&0\\
 0&0&-2&0&1&-2\\
 0&1&-2&0&1&-3\\
 0&0&-1&1&0&-1
 \end{pmatrix}.
 \tag{F.U11}
\]
Take $M_1=A^t$, $M_2=B_1^t$ and
$M_0=(M_1M_2)^{-1}$.  The displayed twists are fixed integral
vectors with $\ell(w_1)=1$, $\ell(w_2)=-1$.
They therefore give free affine actions.  The cusp is
primitive by F.2.16.  Exact Hermite reduction gives
\[
 \operatorname{col}_{\mathbb Z}
 [\,M_0-I\mid M_1(M_0-I)\mid M_2(M_0-I)\,]
       =\langle e_1,\ldots,e_5\rangle.
\]
Van Kampen leaves only $5a=h$, $4b=-h$, $a+b=0$,
whose relation matrix is unimodular.  This proves
$\pi_1=1$ for the specified gluing, including zero cusp
clutching.

For completeness we give the full new integral certificates
needed to reuse the already proved differential calculation
in F.2.6.  The rational exterior invariant and parabolic ranks
are exactly the table there.  The finite cap first homology
groups are $\mathbb Z^2$ and $\mathbb Z^2\oplus\mathbb Z/2$.
On integral invariant two-form bases their suspension periods
divided by the covering orders are, respectively,
$(-1,0,-1)$ and $(0,0,1)$.  The determinants of the invariant
degree-three skew pairings are $5^4$ and $4^4$, while the
degree-two/degree-four pairings have absolute determinants
$5^3$ and $4^3$.  Thus all the finite covering indices used
in degrees two, three and four are one.

For exterior coefficient degrees two, three and four, the
global coboundary Smith form has one factor two and all other
nonzero factors one.  Its torsion class has peripheral
triviality flags $(\mathrm{true},\mathrm{false},\mathrm{true})$,
so the parabolic groups are torsion free, of ranks $0,3,0$.
The degree-one coboundary is primitive and the degree-one
finite covering cokernels are killed by the common invariant
covector and the Chinese remainder theorem, exactly as in
F.2.6.  The global coinvariants in degrees one through four
are each $\mathbb Z/2$.

Write $M_2=\operatorname{diag}(1,D)$.  This time use the
unimodular marking
\[
 P_1=\begin{pmatrix}
 -1&1&-1&0&-1\\
 -1&2&-1&0&-1\\
 -1&0&0&2&-1\\
 1&-1&0&-1&1\\
 0&0&2&-1&1
 \end{pmatrix},\qquad \det P_1=-1.
\]
It satisfies $DP_1=P_1\operatorname{diag}(\text{four-cycle},-1)$.
The cap therefore has the same cellular mapping-torus model
used to prove (F.6).  With the unchanged twist $w_2=-e_0$,
the boundary coordinates remain $\theta=a+4k$, $h=k$;
the suspension representative restricts to $w$ in the base
direction and to $-\operatorname{Norm}_4(w)=0$ in the
auxiliary fibre direction.  Consequently the coefficient in
the torsion differential is again one.

The new vertical representatives in lexicographic exterior
bases are
\[
\begin{split}
 q=2:\ &(0,-1,-1,-1,0),\\
 q=3:\ &(0,0,1,0,0,0,0,0,0,1),\\
 q=4:\ &(0,0,0,0,0,0,0,-1,-1,0).
\end{split}
\]
Each has residue one in its global coinvariant
$\mathbb Z/2$.  Thus all three torsion differentials in (F.6)
are isomorphisms.  The integral Leray table and its free
degree-one differential are therefore exactly those of
F.2.6.  They give $H^2=H^3=0$ and $H^4=\mathbb Z^4$.
Universal coefficients, Hurewicz and Poincare duality give
the stated integral homology and three-connectivity.

The relative almost-complex extension over the pair of pants
in F.2.5 applies: the finite matrices have the same even
real eigenspace dimensions and the cusp is the same marked
Kuhn type.  The purely smooth characteristic-vector and
normal-invariant argument in F.4.1 then gives the stated
oriented diffeomorphism type.  No integrability at the finite
punctures is used for this smooth classification.

Finally, (F.U8) has $(\det H,\det H')=(0,2)$ at $v=1$,
so its irreducible rank-five quotient is not self-dual.
The exact program
\texttt{code/fourfolds\_v1\_homology.py} runs every matrix,
Smith, Hermite, suspension-period, pairing and differential
certificate just displayed.  Its monomial marking was found
inside the stated finite cube and then verified directly;
the proof uses that explicit marking, not search completeness.
\end{proof}

\paragraph{Theorem F.2.19 (Proved: a residual obstruction for every lattice).}
Let a rational rank-six datum have finite linear orders five
and four, a primitive maximal cusp on a stable integral lattice,
and invariant line $L$, with irreducible rank-five quotient $W$
and no rational global coinvariants.  Suppose that at a prime
$p\ne2,5$ the semisimplified reduction of a stable lattice in
$W$ is the sum of the trivial module and an absolutely
irreducible module of dimension four.  Then no stable integral
lattice in that rational datum, with primitive maximal cusp,
and no free finite affine twists can give
$\pi_1=1$ and $H_2=0$ simultaneously.

For every $v\notin\{0,1\}$, the rational family in F.2.16
satisfies this hypothesis at each prime dividing
$v^2-v+1$.  Consequently the exclusion of those parameters
holds for \emph{all} stable lattices and all twists, not just
the lattice in (F.U6).  The two remaining parameters are the
verified survivors of F.2.6 and F.2.18.

\begin{proof}
We first give the lattice-independent mechanism.  For any
stable lattice $V_\mathbb Z$, use its primitive invariant
line $L_\mathbb Z$ and quotient lattice $W_\mathbb Z$.
The semisimplified reduction of a rational representation is
independent of its stable lattice.  One elementary way to see
this is to compare two commensurable lattices by their finite
quotient $T$.  The exact sequence obtained by reduction modulo
$p$ has end terms $T[p]$ and $T/pT$, which have the same
composition factors by the filtration of $T$ by powers of $p$.
Hence the two lattice reductions have equal classes in the
Grothendieck group of finite-length modules.  Chaining such
comparisons proves the assertion for any two lattices.

Thus $W_\mathbb Z/pW_\mathbb Z$ has length two, with one
trivial constituent and one irreducible four-dimensional
constituent.  It has either a trivial submodule or a trivial
quotient (and possibly both).

Suppose first that it has a trivial submodule.  Its inverse
image in $V_\mathbb Z/pV_\mathbb Z$ is an extension of a
trivial module by a trivial module.  Each finite generator
acts on this two-dimensional space by a unipotent matrix.
Since its order is five or four, prime to $p$, it acts as
the identity.  There are therefore at least two global
invariant vectors in degree-one fibre cohomology modulo $p$.
Write their dimension as $a_p\ge2$; this is also the
dimension of the global coinvariants of fibre homology over
$\mathbb F_p$, by the perfect pairing between fibre homology
and cohomology.

The abelianized van Kampen presentation now shows
\[
 b_1(M;\mathbb F_p)\ge a_p-1.
 \tag{F.U12}
\]
Indeed, after the monodromy relations the fibre contributes
$a_p$ dimensions, and the two finite meridians contribute
two more.  The two finite power relations and the single
base relation impose at most three linear relations.
The cusp introduces no further fibre relation in this
abelianization.  On the dual homology lattice its nilpotent
operator is still primitive and maximal, so its angular lattice is
$\ker(M_0-I)=\operatorname{im}(M_0-I)$ by primitivity, and
the latter image already vanishes among the global
monodromy relations.  Its base meridian is the base relation
just counted.  Changing affine twists changes coefficients
of these three relations, not their number.  Thus (F.U12)
contradicts simple connectivity.

Suppose instead that the reduction of $W$ has a trivial
quotient.  Its integral global coinvariants have nonzero
$p$-primary torsion, since they have rational rank zero.
The natural map from the global coinvariants of $V$ to
those of $W$ is an isomorphism.  To check its kernel, a
primitive generator $e$ of $L_\mathbb Z$ lies in
$\ker(M_0-I)=\operatorname{im}(M_0-I)$, and hence is already
zero among global coinvariants.  Therefore $H^2(R^1)$ has
nonzero $p$-primary torsion.  If the total space were simply
connected, the finite-cap argument of F.2.17 would make
every possible torsion source in $H^0(R^2)$ supported only
at two and five.  The invariant part of $\bigwedge^2V$ is
zero by irreducibility of the odd-dimensional quotient.
Consequently the nonzero $p$-primary subgroup survives to
$H^3(M;\mathbb Z)$ and then to
$\operatorname{Tor}H_2(M;\mathbb Z)$, again a contradiction.
This proves the general criterion.

For its application to F.2.16, let $p\mid q=v^2-v+1$.
We have $p\ne2,5$.  In the original cycle lattice of (F.U1),
$Y\mathbf1=2q(-1,1-v)^t$, so $\mathbf1$ gives a trivial
submodule after reduction.  On the four-dimensional quotient
the first finite generator has characteristic polynomial
$\Phi_5$, and $C=B^{-1}$ has polynomial
$(x+1)^2(x^2+1)$.  This quotient is absolutely irreducible.
Here is a characteristic-independent proof over
$\overline{\mathbb F}_p$.

A one-dimensional composition factor would make the
eigenvalues of $A$ and $C$ equal, because $CA^{-1}$ is
unipotent.  An eigenvalue of order dividing both five and
four must be one, which is absent from $\Phi_5$.
If the quotient were reducible, its only remaining possible
composition-factor dimensions would be two and two.
On each factor, the determinants of $A$ and $C$ must agree
and have orders dividing five and four, so both determinants
are one.  The two $C$ eigenvalue pairs must therefore be
$(-1,-1)$ and $(i,-i)$ in some order.  On the former
factor $C=-I$, since its order is prime to $p$.
Then the unipotence of $CA^{-1}=-A^{-1}$ would force the
eigenvalues of $A$ to be $-1$, impossible for $\Phi_5$.
This proves absolute irreducibility and establishes the
required residual composition factors at every $p\mid q$.
When $v\notin\{0,1\}$, the integer $q$ is greater than one
and has such a prime divisor.  The criterion applies.
\end{proof}
